\usepackage{xcolor} % colours for \hl (yellow highlights) and \unclear

\usepackage{float}
\usepackage{graphicx}
\usepackage{xltabular}
\usepackage{enumitem}
\usepackage{makecell}
\usepackage{geometry}
\usepackage[normalem]{ulem}
\usepackage{amsmath,amssymb}
\usepackage{booktabs,array,multirow}
  \geometry{a4paper,inner=3.5cm,outer=2.5cm,top=2.5cm,bottom=2.5cm,includefoot,footskip=1.25cm}
\usepackage[backend=biber,style=vancouver,citestyle=numeric-comp,sortcites=true,sorting=none,maxbibnames=99,giveninits=true,doi=true,url=true]{biblatex}
  \addbibresource{references.bib}
\usepackage{xepersian}

% XePersian requires an explicit Persian text font; fall back like the topic
% fragments do (Amiri → FreeSerif → DejaVu Sans).
\IfFontExistsTF{Amiri}{\settextfont{Amiri}}{%
  \IfFontExistsTF{FreeSerif}{\settextfont{FreeSerif}}{\settextfont{DejaVu Sans}}}
\IfFontExistsTF{DejaVu Serif}{\setlatintextfont{DejaVu Serif}}{%
  \IfFontExistsTF{FreeSerif}{\setlatintextfont{FreeSerif}}{\setlatintextfont{DejaVu Sans}}}



\renewcommand{\contentsname}{فهرست مطالب}
\renewcommand{\listtablename}{فهرست جدول‌ها}
\renewcommand{\listfigurename}{فهرست شکل‌ها و نمودارها}
\renewcommand{\bibname}{فهرست منابع}
\newcommand{\en}[1]{\lr{#1}}
\newcommand{\rng}[2]{\lr{$#1$\textendash$#2$}}
\newcommand{\hl}[1]{\colorbox{yellow!35}{\strut #1}}
\newcommand{\unclear}[1]{\textcolor{red!65!black}{[خوانش نامطمئن: #1]}}
\newcommand{\sourcepage}[2]{\clearpage\section*{صفحهٔ #1}\begingroup\footnotesize\en{#2}\par\endgroup}

\newcommand{\FrontMatterNumbering}{\pagenumbering{harfi}\pagestyle{thesis}}
\newcommand{\MainMatterNumbering}{\pagenumbering{arabic}\pagestyle{thesis}}

% --- notebook build glue -----------------------------------------------------
% \pagestyle{thesis} (used by \FrontMatterNumbering) is not provided by
% XePersian, so define a minimal one: empty header, centred page number.
\makeatletter
\def\ps@thesis{%
  \let\@oddhead\@empty\let\@evenhead\@empty%
  \def\@oddfoot{\hss\footnotesize\thepage\hss}%
  \let\@evenfoot\@oddfoot%
}
\makeatother
% Hook referenced in main.tex before the bibliography chapter (kept as a
% no-op placeholder so the layout stays neutral).
\providecommand{\ChapterStart}{}

% --- per-file InNotebook flags ------------------------------------------------
% index.tex sets a per-file <X>InNotebook flag before \input of each topic file
% whose dual-mode guard checks it (\ifdefined\<X>InNotebook). Define every flag
% used by the current index.tex as a no-op, so the guards take the
% "already inside the notebook" branch instead of re-running their standalone
% preamble mid-document (which ends the book early via \end{document}).
% When index.tex is regenerated with a new topic, add its flag here too.
\providecommand{\AbortionInNotebook}{}
\providecommand{\AGEInNotebook}{}
\providecommand{\AllergicDrugReactionsInNotebook}{}
\providecommand{\AllergicRhinitisInNotebook}{}
\providecommand{\AnaphylaxisInNotebook}{}
\providecommand{\AnemiaHyperbilirubinemiaInNotebook}{}
\providecommand{\CAPInNotebook}{}
\providecommand{\CarbohydrateDisordersInNotebook}{}
\providecommand{\CeliacInNotebook}{}
\providecommand{\ComplementInNotebook}{}
\providecommand{\CommonColdInNotebook}{}
\providecommand{\DKAInNotebook}{}
\providecommand{\DisturbanceInNotebook}{}
\providecommand{\ElectrolyteDisordersInNotebook}{}
\providecommand{\EmergencyContraceptionInNotebook}{}
\providecommand{\FWLSInNotebook}{}
\providecommand{\FUOInNotebook}{}
\providecommand{\FeverInNotebook}{}
\providecommand{\FeverOlderInNotebook}{}
\providecommand{\FeverRashInNotebook}{}
\providecommand{\FetalDiseasesInNotebook}{}
\providecommand{\FetalGrowthInNotebook}{}
\providecommand{\GDMInNotebook}{}
\providecommand{\GERDInNotebook}{}
\providecommand{\HTNInNotebook}{}
\providecommand{\HypernatremiaInNotebook}{}
\providecommand{\IDMInNotebook}{}
\providecommand{\IntussusceptionInNotebook}{}
\providecommand{\LBPInNotebook}{}
\providecommand{\LeukemiaInNotebook}{}
\providecommand{\LymphadenopathyInNotebook}{}
\providecommand{\LymphocyteDisordersInNotebook}{}
\providecommand{\MaternalDiseasesInNotebook}{}
\providecommand{\NECInNotebook}{}
\providecommand{\NeonatalRespiratoryInNotebook}{}
\providecommand{\NeonatalSeizuresInNotebook}{}
\providecommand{\NeonateEvaluationInNotebook}{}
\providecommand{\RDSInNotebook}{}
\providecommand{\SSInNotebook}{}
\providecommand{\STIInNotebook}{}
\providecommand{\SplenomegalyInNotebook}{}

% \marktext marks a key phrase transcribed from highlighted/boxed handwriting in
% the source pages (used by neurology/headache; body-only file, so the parent
% must supply it). Same visual convention as the sibling topic files.
\providecommand{\marktext}[1]{\begingroup\setlength{\fboxsep}{1pt}\colorbox{yellow!40}{\strut #1}\endgroup}

% main.tex uses \phantomsection before the bibliography, but hyperref is not
% loaded in this XePersian setup; provide a no-op so the anchor line is harmless.
\providecommand{\phantomsection}{}
